% Options for packages loaded elsewhere
\PassOptionsToPackage{unicode}{hyperref}
\PassOptionsToPackage{hyphens}{url}
\documentclass[
]{article}
\usepackage{xcolor}
\usepackage[margin=1in]{geometry}
\usepackage{amsmath,amssymb}
\setcounter{secnumdepth}{-\maxdimen} % remove section numbering
\usepackage{iftex}
\ifPDFTeX
  \usepackage[T1]{fontenc}
  \usepackage[utf8]{inputenc}
  \usepackage{textcomp} % provide euro and other symbols
\else % if luatex or xetex
  \usepackage{unicode-math} % this also loads fontspec
  \defaultfontfeatures{Scale=MatchLowercase}
  \defaultfontfeatures[\rmfamily]{Ligatures=TeX,Scale=1}
\fi
\usepackage{lmodern}
\ifPDFTeX\else
  % xetex/luatex font selection
\fi
% Use upquote if available, for straight quotes in verbatim environments
\IfFileExists{upquote.sty}{\usepackage{upquote}}{}
\IfFileExists{microtype.sty}{% use microtype if available
  \usepackage[]{microtype}
  \UseMicrotypeSet[protrusion]{basicmath} % disable protrusion for tt fonts
}{}
\makeatletter
\@ifundefined{KOMAClassName}{% if non-KOMA class
  \IfFileExists{parskip.sty}{%
    \usepackage{parskip}
  }{% else
    \setlength{\parindent}{0pt}
    \setlength{\parskip}{6pt plus 2pt minus 1pt}}
}{% if KOMA class
  \KOMAoptions{parskip=half}}
\makeatother
\usepackage{longtable,booktabs,array}
\newcounter{none} % for unnumbered tables
\usepackage{calc} % for calculating minipage widths
% Correct order of tables after \paragraph or \subparagraph
\usepackage{etoolbox}
\makeatletter
\patchcmd\longtable{\par}{\if@noskipsec\mbox{}\fi\par}{}{}
\makeatother
% Allow footnotes in longtable head/foot
\IfFileExists{footnotehyper.sty}{\usepackage{footnotehyper}}{\usepackage{footnote}}
\makesavenoteenv{longtable}
\setlength{\emergencystretch}{3em} % prevent overfull lines
\providecommand{\tightlist}{%
  \setlength{\itemsep}{0pt}\setlength{\parskip}{0pt}}
\usepackage{mathtools}
\usepackage{enumitem}
\usepackage{microtype}
\usepackage{booktabs}
\usepackage{array}
\usepackage{amssymb}
\usepackage{newunicodechar}
% Greek
\newunicodechar{α}{\ensuremath{\alpha}}
\newunicodechar{δ}{\ensuremath{\delta}}
\newunicodechar{η}{\ensuremath{\eta}}
\newunicodechar{λ}{\ensuremath{\lambda}}
\newunicodechar{μ}{\ensuremath{\mu}}
\newunicodechar{ξ}{\ensuremath{\xi}}
\newunicodechar{π}{\ensuremath{\pi}}
% superscripts / subscripts (text-native so they never demand math mode, even in \texttt)
\newunicodechar{¹}{\textsuperscript{1}}
\newunicodechar{²}{\textsuperscript{2}}
\newunicodechar{³}{\textsuperscript{3}}
\newunicodechar{⁴}{\textsuperscript{4}}
\newunicodechar{⁵}{\textsuperscript{5}}
\newunicodechar{⁶}{\textsuperscript{6}}
\newunicodechar{⁷}{\textsuperscript{7}}
\newunicodechar{⁻}{\textsuperscript{$-$}}
\newunicodechar{₀}{\textsubscript{0}}
% math operators / relations
\newunicodechar{√}{\ensuremath{\surd}}
\newunicodechar{−}{\ensuremath{-}}
\newunicodechar{×}{\ensuremath{\times}}
\newunicodechar{≈}{\ensuremath{\approx}}
\newunicodechar{∼}{\ensuremath{\sim}}
\newunicodechar{ℓ}{\ensuremath{\ell}}
\newunicodechar{∂}{\ensuremath{\partial}}
\newunicodechar{∇}{\ensuremath{\nabla}}
\newunicodechar{∈}{\ensuremath{\in}}
\newunicodechar{∝}{\ensuremath{\propto}}
\newunicodechar{∫}{\ensuremath{\int}}
\newunicodechar{≠}{\ensuremath{\neq}}
\newunicodechar{≤}{\ensuremath{\leq}}
\newunicodechar{≥}{\ensuremath{\geq}}
\newunicodechar{≲}{\ensuremath{\lesssim}}
\usepackage{bookmark}
\IfFileExists{xurl.sty}{\usepackage{xurl}}{} % add URL line breaks if available
\urlstyle{same}
\hypersetup{
  hidelinks,
  pdfcreator={LaTeX via pandoc}}

\title{ED Gravity Is a Healthy One-Parameter Khronometric Deformation of General Relativity}
\author{Allen Proxmire}
\date{July 2026}

\begin{document}
\maketitle

\textbf{Series:} Event Density (ED) Generative Papers, gravity arc
(curvature emergence) \textbf{Status:} Publication draft (structural
consolidation). Consolidates the \emph{dynamical} half of the ED gravity
arc: why ED's emergent gravity is a preferred-frame (khronometric)
theory, what that fixes (GR-like tensor radiation, a shared light cone,
a scalar that supplies MOND), the form of its effective action, and the
computation that bounds its one distinctive parameter λ to the healthy
low branch {[}0, 1/3). Standalone; cold-reader accessible. Companion to
the kinematics paper (Relational Gravity: the Emergent Metric) and to
Paper\_027 (Newton's G). \textbf{Repository target:}
\texttt{physics-papers/gravity/} (ED-Generative)

\begin{center}\rule{0.5\linewidth}{0.5pt}\end{center}

\subsection{\texorpdfstring{Preamble: What This Paper Does NOT Claim
\emph{(written first per QC discipline; abstract reconciled against
this)}}{Preamble: What This Paper Does NOT Claim (written first per QC discipline; abstract reconciled against this)}}\label{preamble-what-this-paper-does-not-claim-written-first-per-qc-discipline-abstract-reconciled-against-this}

\begin{enumerate}
\def\labelenumi{\arabic{enumi}.}
\item
  \textbf{This is not a derivation of exact General Relativity.} ED
  gravity is a \emph{khronometric} (Lorentz-violating, preferred-frame)
  theory in the Einstein-aether / Hořava class. It shares GR's spin-2
  tensor graviton (so its tensor radiation is GR-like), but it carries
  an extra scalar (khronon) mode and its metric kinetic structure
  differs from GR's in the trace sector. The Einstein equation is not
  derived. ED \emph{reduces to} GR at high acceleration and gives MOND
  at low acceleration; it is not GR.
\item
  \textbf{The khronometric identification is an account-tier structural
  argument, not a from-nothing derivation.} The claim ``ED has a
  preferred frame (the arrow, P11), therefore its emergent gravity is
  Lorentz-violating, therefore it is the khronometric class'' is a
  structural identification with a known, standard,
  observationally-viable class of gravity theories. It is grounded in an
  established ED fact (the arrow is a genuine preferred frame), but the
  passage from ``preferred frame'' to ``this specific effective action''
  uses the standard effective-field-theory logic for such theories, not
  a bespoke ED derivation of every coupling.
\item
  \textbf{λ is bounded and signed, not pinned to a decimal.} The result
  \texttt{λ\ ∈\ {[}0,\ 1/3)} (natural value 0, and \texttt{λ\ ≠\ 1}
  given ED's positive-energy substrate inertia) is a \emph{structural
  argument}: the exact decomposition and the three reference values are
  derived, but ED's placement within the interval rests on the
  association that the substrate's coherence / participation-velocity
  channel supplies the kinetic term (see \#4). The exact point in {[}0,
  1/3), which would fix the precise scalar-mode speed and the MOND
  interpolation constant, is not computed here; it needs the
  mode-resolved conservation inertia and is stated as open.
\item
  \textbf{The coherence-to-kinetic map is an association, honestly
  flagged.} Identifying the effective action's kinetic term with ED's
  coherence sector (the reversible participation channel) is organized
  by a non-canonical source (an older ED cosmology manuscript) and the
  author's ``Relation-Boundary-Gradient'' reading of the P12 landscape.
  It is a useful pointer that locates \emph{which} part of the substrate
  supplies the inertia; it is not a canonical primitive statement. The
  canonical source is Paper\_087. The one piece of the map that is
  \emph{verified} (the gradient sector → spatial curvature) is marked as
  such.
\item
  \textbf{c\_T = c is a confirmed consistency, not a tuned fit and not a
  forced prediction.} A shared photon/graviton light cone (the
  gravity-sector condition c\_13 = 0) is a \emph{separate} condition
  from the matter-sector universality that evades the dipole; ED's
  minimal setup makes the shared cone natural, and GW170817 confirms it
  to 10⁻¹⁵. This is offered as a distinctive consistency (a condition ED
  finds natural, confirmed by observation), not as a free parameter
  adjusted to data, and not as a coupling ED derives from first
  principles (the action that would fix c\_13 is unbuilt). Its
  load-bearing assumption (that ED's universality is exact enough to
  hold c\_T = c to that precision through the radiative sector) is
  stated.
\item
  \textbf{Values are inherited.} Newton's G, the substrate scale ℓ\_P,
  and the MOND acceleration a₀ are inherited from measurement or from
  Paper\_027; this paper derives forms and bounds, not those numbers.
\item
  \textbf{No new primitive is introduced.} The preferred frame is P11
  (the arrow), already primitive.
\end{enumerate}

\begin{center}\rule{0.5\linewidth}{0.5pt}\end{center}

\subsection{Abstract}\label{abstract}

The companion kinematics paper established that a metric emerges from
ED's participation graph and is fixed to the weak-field spatial factor
\texttt{g\ \textasciitilde{}\ 1/b} in 3D, with a Gauss's-law Newtonian
limit and a MOND-resembling nonlinearity, but a kinematic (acoustic)
metric appears not to radiate, which binary-pulsar decay and
gravitational-wave inspirals demand. This paper resolves the
\emph{dynamical} sector. The resolution turns on one established ED
fact: \textbf{ED has a preferred frame, the arrow (P11).} An emergent
gravity with a preferred frame is Lorentz-violating, and the effective
field theory of such a gravity is the \textbf{khronometric /
Einstein-aether / Hořava class} with the aether four-vector equal to the
arrow. This is not a scalar theory: it carries the \textbf{spin-2 tensor
graviton}, so its tensor radiation is GR-like and matches binary-pulsar
orbital decay and LIGO's tensor polarization, plus an extra scalar
(khronon) mode. A probe confirms the emergent metric is \emph{not}
conformal-only: a non-spherical source develops a genuine traceless
spin-2 part (0.000 for a point mass, nonzero and growing with anisotropy
for a quadrupolar source), so the point-mass conformal
\texttt{g\ \textasciitilde{}\ 1/b} was a special case and the
``scalar-only'' reading that motivated an earlier ``1/6 falsified''
luminosity is retracted (the static probe shows the channel is
\emph{present}; its \emph{propagation} as a GR-like graviton follows
from the class, not the probe). ED's minimal, universal setup then
selects the viable corner, via three \emph{distinct} conditions on
\emph{different sectors} (not one mechanism): its universal coupling
(equivalence principle, mass = universal bandwidth-influence
\texttt{Q\ ∝\ M}) structurally kills the scalar \emph{dipole} radiation
(the strongest scalar-gravity pulsar test, the one genuinely
ED-structural leg); a shared graviton/photon light cone (the separate
gravity-sector condition \texttt{c\_13\ =\ 0}, giving \textbf{c\_T = c},
consistent with GW170817's 10⁻¹⁵ bound); and small preferred-frame PPN
couplings (a third condition, \texttt{c\_14\ →\ 0}). The dipole evasion
is an ED-structural consequence; the shared cone and small PPN are
structural conditions ED's minimality makes natural but does not derive
(open), so \texttt{c\_T\ =\ c} is a confirmed \emph{consistency}, not a
forced prediction. The effective action takes the khronometric form
\(S = (1/16\pi G)\int N\sqrt{h}\,[K_{ij}K^{ij} - \lambda K^2 + \xi\,{}^3R + \eta\,a^2]\),
and ED's structural conditions give \texttt{ξ\ =\ 1} (shared light cone)
and \texttt{η\ →\ 0} (equivalence principle), leaving \texttt{λ}, the
coefficient of \texttt{K²}, as the \emph{single} distinctive parameter.
We then compute \texttt{λ}. The kinetic term decomposes exactly (machine
precision) into a shear (spin-2) part and a compression (spin-0, trace)
part, and \texttt{λ} is nothing but their kinetic-weight ratio, with
three derived landmarks: \texttt{λ\ =\ 0} (compression a normal
positive-energy mode, the minimal isotropic substrate),
\texttt{λ\ =\ 1/3} (compression non-dynamical), and \texttt{λ\ =\ 1}
(GR, where compression is a negative-energy conformal ghost, a value
\emph{forced} by full four-dimensional diffeomorphism invariance, which
a khronometric substrate does not possess). ED's coherence /
participation-velocity inertia gives compression and shear equal
positive weight, hence the natural value \texttt{λ\ =\ 0}; bandwidth
conservation ties the compression mode to the current, which can only
suppress its weight toward \texttt{1/3}, never past it into GR's ghost.
So \textbf{\texttt{λ\ ∈\ {[}0,\ 1/3)}, natural value 0, \texttt{λ\ ≠\ 1}
(given ED's positive-energy substrate inertia), the healthy low branch}
on which the extra scalar is a well-behaved positive-energy field (the
standard khronometric stability window is \texttt{λ\ \textless{}\ 1/3}).
The clean unconditional part is that ED is \emph{not forced} to GR's
\texttt{λ\ =\ 1} (a khronometric substrate lacks the 4D diffeomorphism
invariance that forces it); the natural value 0 and the
\texttt{{[}0,\ 1/3)} placement are a structural argument resting on a
flagged association (that the substrate's coherence channel supplies the
kinetic term). Because the shear (tensor) weight is 1 in ED and GR
alike, \texttt{λ} never touches the tensor sector: ED radiates a GR-like
graviton regardless, and \texttt{λ} sets only the scalar (khronon)
sector, the candidate home for MOND. The honest verdict: \textbf{ED
gravity is a healthy one-parameter khronometric deformation of GR} with
the arrow as its preferred frame, GR-like tensor radiation, a shared
light cone (\texttt{c\_T\ =\ c}, consistent with GW170817), a healthy
scalar in the sector where galactic MOND would live (that it reproduces
MOND is inherited/open), and its one distinctive parameter bounded to
\texttt{{[}0,\ 1/3)}. Open: the exact point in \texttt{{[}0,\ 1/3)}
(needing the mode-resolved conservation inertia), which fixes the
precise scalar-mode speed and MOND interpolation constant.

\begin{center}\rule{0.5\linewidth}{0.5pt}\end{center}

\subsection{1. Introduction: the Metric Emerges; Does It Have Dynamics,
and of What
Class?}\label{introduction-the-metric-emerges-does-it-have-dynamics-and-of-what-class}

The companion kinematics paper answered where ED's geometry comes from:
a metric emerges from bandwidth-connectivity, is fixed to
\texttt{g\ \textasciitilde{}\ 1/b} in 3D by the holographic
channel-count, obeys a Gauss's-law Newtonian field equation, is
relational (lives on graph adjacency), and has a MOND-resembling
nonlinearity because the metric is a kinematic read-out of the bandwidth
field. That is the \emph{static and linear} story, and it left one sharp
tension unresolved: a kinematic (acoustic) metric does not obviously
radiate energy as self-propagating spin-2 waves, yet binary-pulsar
orbital decay (matched to GR at sub-percent) and gravitational-wave
inspirals are \emph{weak-field, radiative} confirmations of
dynamical-metric behaviour. This is not deferrable to a strong-field
hedge; it is a live weak-field question.

This paper resolves it, and the resolution is not a patch but a
reclassification of what \emph{kind} of gravity ED is. One established
fact does the work: \textbf{ED has a preferred frame, the arrow (P11),
the commitment/irreversibility primitive that orders events.} A gravity
theory built on an emergent metric that inherits a preferred frame is
Lorentz-violating, and Lorentz-violating gravity has a well-developed
effective field theory: the Einstein-aether / Hořava / khronometric
class, with a dynamical unit timelike vector (the aether) picking out
the preferred foliation. Identifying that aether with the arrow places
ED in this class. The consequences are specific and largely favourable:
such theories carry GR's spin-2 tensor graviton (so they radiate like
GR), they admit a shared light cone (so \texttt{c\_T\ =\ c} is natural),
and they carry one extra scalar mode whose low-acceleration behaviour is
the natural home for MOND.

Section 3 makes the preferred-frame identification. Section 4 shows the
emergent metric carries a genuine spin-2 channel (so ED is not a scalar
theory). Section 5 shows how ED's minimality places it in the viable
radiative corner via three distinct conditions (dipole evasion,
\texttt{c\_T\ =\ c}, PPN suppression), only the first of which is
ED-structural. Section 6 gives the effective action and reduces it to
the single parameter \texttt{λ}. Section 7 bounds \texttt{λ} to the
healthy low branch \texttt{{[}0,\ 1/3)}. Section 8 gives the honest
substrate map behind that bound. Sections 9-11 are limitations,
falsifiers, and the position statement.

\subsection{2. Load-Bearing Step Audit}\label{load-bearing-step-audit}

{\def\LTcaptype{none} % do not increment counter
\begin{longtable}[]{@{}
  >{\raggedright\arraybackslash}p{(\linewidth - 4\tabcolsep) * \real{0.3333}}
  >{\raggedright\arraybackslash}p{(\linewidth - 4\tabcolsep) * \real{0.3333}}
  >{\raggedright\arraybackslash}p{(\linewidth - 4\tabcolsep) * \real{0.3333}}@{}}
\toprule\noalign{}
\begin{minipage}[b]{\linewidth}\raggedright
Step
\end{minipage} & \begin{minipage}[b]{\linewidth}\raggedright
Status
\end{minipage} & \begin{minipage}[b]{\linewidth}\raggedright
Source / justification
\end{minipage} \\
\midrule\noalign{}
\endhead
\bottomrule\noalign{}
\endlastfoot
ED has a preferred frame (the arrow, P11) & P / established &
Paper\_087; the arrow is a genuine preferred frame (irreversibility,
event ordering) \\
Emergent gravity with a preferred frame is Lorentz-violating →
khronometric class & \textbf{grounded / identification} & §3; standard
EFT of preferred-frame gravity (Einstein-aether / Hořava), aether =
arrow \\
Primitive arrow fixes the foliation; the \emph{dynamical} khronon
(b-sourced) propagates the scalar & \textbf{structural assumption} & §3;
a rigid arrow alone would carry no scalar; the emergent b-fluctuation is
assumed dynamical \\
The emergent metric carries a genuine spin-2 channel (non-spherical
source) & \textbf{present (static probe); propagation from the class} &
§4; probe: traceless part 0.000 (point mass), nonzero (quadrupolar);
static, not a radiative proof \\
Khronometric class radiates a GR-like spin-2 graviton &
\textbf{inherited (standard result)} & §4-5; the class carries the
tensor graviton; tensor radiation GR-like \\
Scalar \emph{dipole} radiation evaded by universal coupling (equivalence
principle) & \textbf{ED-structural} & §5; mass = universal
\texttt{Q\ ∝\ M} → scalar dipole ∝ mass dipole, \texttt{D̈\ =\ 0};
probe-confirmed \\
\texttt{c\_T\ =\ c} (shared graviton/photon cone, \texttt{c\_13\ =\ 0})
& \textbf{structural condition (gravity sector), confirmed consistency}
& §5; a \emph{separate} condition from matter universality; GW170817
confirms 10⁻¹⁵; not derived \\
Preferred-frame PPN parameters suppressed (\texttt{c\_14\ →\ 0}) &
\textbf{structural condition, open} & §5-6; a \emph{third} condition;
\texttt{α\_1\ =\ −4\ c\_14} on \texttt{c\_13=0};
\texttt{c\_14}-smallness natural but not derived \\
Effective action form (khronometric) + reduction to one parameter λ &
\textbf{grounded / account} & §6; \texttt{ξ=1} (light cone),
\texttt{η→0} (equivalence principle), \texttt{λ} = the one distinctive
coupling; \texttt{λ\ =\ 1+c\_2} on \texttt{c\_13=0} \\
Kinetic decomposition
\texttt{K\_ij\ K\^{}ij\ −\ λK²\ =\ A²\ +\ (1/3−λ)K²}; the three λ
landmarks & \textbf{derived} & §7; exact algebra (trace/traceless
orthogonality) + standard DeWitt / Hořava facts \\
GR's λ=1 forced by 4D diffeomorphism invariance; ED \emph{not forced} to
it & \textbf{derived (non-circular)} & §7; the conformal-ghost value is
a covariance artifact ED's substrate lacks; clean unconditional part \\
λ = 0 natural (isotropic participation-velocity inertia) &
\textbf{structural argument (association)} & §7-8; Coh → kinetic is a
non-canonical association (flagged) \\
λ ∈ {[}0, 1/3), healthy low branch, λ ≠ 1 & \textbf{structural argument
(on the positive-energy premise)} & §7; natural 0 (association) +
\texttt{1/3} = fully-constrained endpoint + positive-energy premise;
interior open \\
Grad → spatial curvature (³R) & \textbf{verified} & §8; probe:
\texttt{³R\ =\ 2∇²b\ −\ (5/2)(∇b)²/b}, the gradient content of
\texttt{b} \\
Exact point of λ in {[}0, 1/3); precise scalar speed and MOND constant &
\textbf{OPEN} & §7-8; needs the mode-resolved conservation inertia \\
Exact GR (Einstein equation) & \textbf{NOT CLAIMED} & preamble 1;
khronometric, not GR \\
\texttt{G}, \texttt{ℓ\_P}, \texttt{a₀} values & \textbf{inherited} &
preamble 6 \\
\end{longtable}
}

Every step is P, established, inherited, demonstrated, structural,
derived, grounded, verified, a flagged association, or explicitly open /
not-claimed. No step claims exact GR or a computed decimal for λ.

\subsection{3. The Arrow Is a Preferred Frame, so ED Gravity Is
Khronometric}\label{the-arrow-is-a-preferred-frame-so-ed-gravity-is-khronometric}

ED's arrow (P11) is the commitment/irreversibility primitive: it orders
events, distinguishes past from future, and is not a coarse-graining
artifact but a substrate primitive (the reduction program isolates it as
the keystone). A preferred temporal ordering is exactly a
\emph{preferred frame}: a physically distinguished foliation of the
emergent spacetime into ``space at a time.'' Relativity's local Lorentz
invariance has no such distinguished foliation; ED's arrow does. So ED's
emergent gravity is \textbf{Lorentz-violating at the fundamental level},
with local Lorentz invariance an emergent, approximate symmetry of the
low-energy sector rather than an exact one.

Lorentz-violating gravity is not exotic; it has a mature effective field
theory. Represent the preferred frame by a dynamical unit timelike
vector field \texttt{u\^{}μ} (the ``aether''), normalized
\texttt{u\^{}μ\ u\_μ\ =\ 1}. The most general two-derivative action
coupling \texttt{u} to the metric is Einstein-aether theory; when
\texttt{u} is further required to be hypersurface-orthogonal (to define
a global foliation, a ``khronon'' scalar time function \texttt{T} with
\texttt{u\_μ\ ∝\ ∂\_μ\ T}), it is khronometric theory, the low-energy
limit of Hořava gravity. \textbf{ED's arrow is precisely such a global
time function}: it is a universal ordering, hypersurface-orthogonal by
construction (all observers share the substrate's commitment order). So
the natural identification is

\[ u_\mu \;=\; \text{the arrow (P11)}, \qquad \text{ED gravity} \;\in\; \text{the khronometric class.} \]

This identification is account-tier (preamble 2): it is grounded in the
established fact that the arrow is a genuine preferred frame, and it
uses the standard EFT of preferred-frame gravity rather than a bespoke
ED derivation. But it is highly constraining, because the khronometric
class has a fixed field content and a fixed set of observational
handles, which the rest of this paper works through. The single most
important consequence is immediate: \textbf{the khronometric class is
not scalar gravity.} It carries the same spin-2 tensor graviton as GR,
plus one extra scalar. So the ``ED gravity is a scalar theory, therefore
falsified by pulsars'' worry is dissolved at the level of field content,
before any calculation.

\textbf{Two things must be kept distinct here, and their distinction is
load-bearing.} The \emph{primitive} arrow (P11) is a rigid,
non-dynamical fact: it fixes that a preferred foliation \emph{exists}
and its orientation (the direction of commitment). By itself a rigid
background aether would carry no propagating scalar, and then there
would be no khronon mode and no MOND sector: a rigid foliation is not
enough. What supplies the propagating scalar is the \emph{dynamical}
khronon, the emergent time-function whose \emph{fluctuations} are
carried by the bandwidth field \texttt{b} (the foliation's local
rate/shape can vary as \texttt{b} varies, exactly the \texttt{∂\_t\ b}
compression mode of §7). So the identification is really two-layered:
the primitive arrow fixes the \emph{existence and orientation} of the
preferred frame (the background), and the emergent, \texttt{b}-sourced
khronon fluctuation is the \emph{dynamical} field that propagates the
scalar mode and hosts MOND. This paper treats the dynamical khronon as
the fluctuation sector of the emergent foliation carried by \texttt{b};
that this emergent khronon is genuinely dynamical (rather than the arrow
being purely rigid, which would remove the scalar entirely) is itself a
structural assumption of the account, flagged here and revisited in §7
(where the compression mode \texttt{∂\_t\ b} \emph{is} that dynamical
content) and the falsifiers (§10.1).

\subsection{4. The Emergent Metric Carries a Genuine Spin-2
Channel}\label{the-emergent-metric-carries-a-genuine-spin-2-channel}

The claim that ED carries a tensor graviton must be checked against the
emergent metric itself, not merely asserted from the class. The
companion paper's \texttt{g\ \textasciitilde{}\ 1/b} is conformal
(\texttt{h\_ij\ ∝\ δ\_ij}, pure trace, spin-0), which \emph{looks} like
a scalar theory. But that form was derived only for a \emph{static,
spherically-symmetric} source, where isotropy forces the conformal form.
A non-spherical source is different: bandwidth then \emph{flows} (a
current, not just a density), and the emergent metric can carry a
traceless part.

A probe (\texttt{spin2\_dof\_probe.py}) tests this directly. It builds
the emergent metric perturbation for sources of varying anisotropy and
extracts the traceless (spin-2) component:

\begin{itemize}
\tightlist
\item
  \textbf{Point / spherical source:} traceless part \texttt{=\ 0.000}
  (the conformal special case that motivated the scalar-only
  misreading).
\item
  \textbf{Quadrupolar (binary-like) source:} traceless part nonzero
  (\texttt{+0.005} at the tested anisotropy), growing with source
  anisotropy.
\end{itemize}

So the emergent metric is \emph{not} conformal-only: the traceless
(spin-2) channel is \emph{kinematically present} for a non-spherical
source. This is necessary, not sufficient, for a \emph{radiative} spin-2
degree of freedom: the probe is static (it shows the channel exists),
and the \emph{propagation} of that channel as a GR-like graviton follows
from the khronometric class identification (§3), not from the probe.
What the probe does establish is enough to retire the ``scalar-only''
reading (a binary's time-varying quadrupole sources a time-varying
traceless part, so the metric is not a pure scalar), and with it the
earlier ``1/6 scalar luminosity, falsified'' conclusion, which had
dropped the tensor sector entirely (see the companion paper's §6
retraction).

\subsection{5. ED's Minimal, Universal Setup Selects the Viable Corner
(Three Conditions, One
Naturalness)}\label{eds-minimal-universal-setup-selects-the-viable-corner-three-conditions-one-naturalness}

The khronometric class is viable only in a corner of its parameter
space, bounded by three sharp constraints: no excess (dipole) radiation,
\texttt{c\_T\ =\ c} (GW170817), and small preferred-frame effects (PPN).
These are \emph{three distinct conditions on different sectors} of the
theory, and it is essential not to conflate them: dipole radiation is a
\emph{matter-sector} question, the graviton light cone is a \emph{pure
gravity+aether} coupling (\texttt{c\_13}), and the PPN parameters turn
on a \emph{further} aether coupling (\texttt{c\_14}), and \texttt{c\_13}
and \texttt{c\_14} are independent points in coupling space. ED does
\textbf{not} derive all three from a single coupling. What ED offers is
a \emph{naturalness}: its setup is universal and minimal (the emergent
metric is one universal read-out of the bandwidth field \texttt{b}, with
no independent aether-matter couplings tuned against each other), which
makes the minimal corner where all three conditions hold the natural
one. Taking the three in turn, and tiering each honestly:

\begin{itemize}
\item
  \textbf{Dipole radiation evaded (matter sector, ED-structural).}
  Scalar gravity generically radiates at \emph{dipole} order, the
  strongest pulsar killer of scalar-tensor theories. But in ED a mass is
  a \emph{universal} bandwidth-influence, \texttt{Q\ ∝\ M} (all bodies
  fall alike, because the metric is one universal shadow of \texttt{b},
  the companion paper's kinematic-metric result), so the scalar dipole
  moment \texttt{D\ =\ Σ\ Q\_i\ x\_i\ =\ (Q/M)\ Σ\ M\_i\ x\_i} is
  proportional to the mass dipole, whose second time-derivative is
  \texttt{d(total\ momentum)/dt\ =\ 0} for an isolated binary. So there
  is no scalar dipole radiation; the leading scalar multipole is the
  quadrupole, as in GR (probe: dipole/quadrupole power
  \texttt{\textasciitilde{}10⁻²²} for universal coupling, reappearing at
  \texttt{\textasciitilde{}10⁻³} if universality is broken). This is the
  one leg that is a genuine ED-\emph{structural} consequence: ED passes
  the strongest scalar-gravity pulsar test because of its equivalence
  principle.
\item
  \textbf{\texttt{c\_T\ =\ c} (tensor-cone condition
  \texttt{c\_13\ =\ 0}; a structural condition ED finds natural,
  confirmed by GW170817).} The tensor graviton's speed equals light's
  when the aether tensor-sector coupling \texttt{c\_13\ =\ 0}, i.e.~the
  graviton and photon share a light cone. This is a
  \emph{gravity-sector} condition, and it does \textbf{not} follow from
  the matter-sector universality that evades the dipole; the two are
  separate. ED's minimal/universal setup (no relative tilt between the
  matter and gravity cones) makes the shared cone natural, and
  GW170817's near-simultaneous gamma-ray and gravitational-wave arrival
  confirmed
  \texttt{\textbar{}c\_T\ −\ c\textbar{}/c\ \textless{}\ 10⁻¹⁵}. So this
  is a \emph{structural condition ED finds natural, confirmed by
  observation}, a consistency, not a coupling ED derives from first
  principles (the coarse-grained action that would fix \texttt{c\_13} is
  unbuilt). (Load-bearing assumption, preamble 5: ED's universality is
  exact enough to hold \texttt{c\_T\ =\ c} to that precision through the
  radiative sector.)
\item
  \textbf{PPN suppressed (\texttt{c\_14\ →\ 0}; a third structural
  condition, open).} The preferred-frame post-Newtonian parameters
  \texttt{α\_1,\ α\_2} encode observable effects of the aether's motion
  relative to matter; on the \texttt{c\_13\ =\ 0} slice
  \texttt{α\_1\ =\ −4\ c\_14} (and \texttt{α\_2} has a related but more
  complex \texttt{c\_14}-dependence). The solar-system and pulsar bounds
  \texttt{\textbar{}α\_1\textbar{}\ ≲\ 10⁻⁴},
  \texttt{\textbar{}α\_2\textbar{}\ ≲\ 10⁻⁷} then require
  \texttt{c\_14\ →\ 0}, a \emph{third} condition, in the aether sector
  again. ED's minimality makes small \texttt{c\_14} natural, but ED does
  \textbf{not} derive it: the smallness of \texttt{c\_14} is open (§6).
\end{itemize}

So ED's universal, minimal setup makes the viable corner
(\texttt{Q\ ∝\ M} universal, \texttt{c\_13\ =\ 0}, \texttt{c\_14\ →\ 0})
the natural one, and the sharpest observation (\texttt{c\_T\ =\ c},
GW170817) is consistent with it. This is a \emph{placement by
naturalness}, account-tier: only the dipole evasion is an ED-structural
consequence; the shared cone and small PPN are separate structural
conditions ED finds natural but does not derive. It is not a derivation
of the three couplings from one mechanism, and not a computation of
exactly where in that corner ED sits, which needs the effective action.

\subsection{6. The Effective Action, Reduced to One Distinctive
Parameter}\label{the-effective-action-reduced-to-one-distinctive-parameter}

The khronometric effective action, in the preferred foliation with lapse
\texttt{N}, spatial metric \texttt{h\_ij}, extrinsic curvature
\texttt{K\_ij}, spatial scalar curvature \texttt{³R}, and aether
acceleration \texttt{a\_i\ =\ ∂\_i\ \textbackslash{}ln\ N}, is

\[ S \;=\; \frac{1}{16\pi G}\int dt\, d^3x\; N\sqrt{h}\;\Big[\, K_{ij}K^{ij} \;-\; \lambda\, K^2 \;+\; \xi\, {}^3R \;+\; \eta\, a_i a^i \,\Big], \]

with the aether \texttt{u} equal to the arrow. GR (in this foliation) is
the special point \texttt{(λ,\ ξ,\ η)\ =\ (1,\ 1,\ 0)}. ED's structural
conditions from §5 point to (they do not \emph{derive}) two of the three
couplings:

\begin{itemize}
\tightlist
\item
  \textbf{\texttt{ξ\ =\ 1}} from the shared light cone
  (\texttt{c\_13\ =\ 0}, \texttt{c\_T\ =\ c}, consistent with GW170817):
  the tensor sector matches GR. A structural condition ED finds natural,
  confirmed by observation, not a derived coupling.
\item
  \textbf{\texttt{η\ →\ 0}} (equivalently \texttt{c\_14\ →\ 0}): the
  aether-acceleration coupling, which sources preferred-frame effects,
  is driven small by the PPN bounds. ED's minimality makes this natural,
  but its smallness is open (§5), not derived.
\end{itemize}

That leaves \textbf{\texttt{λ}}, the coefficient of \texttt{K²}, as the
\emph{single} distinctive parameter, the one that measures ED's
departure from GR. \texttt{λ} lives entirely in the scalar (khronon)
sector: it sets whether and how the extra scalar mode is dynamical,
hence the character of the galactic MOND sector, while the tensor sector
(set by \texttt{ξ\ =\ 1}) radiates like GR regardless. So \textbf{ED
gravity is, to the extent these two structural conditions hold, a
one-parameter khronometric deformation of GR, and that parameter is
\texttt{λ}.} Bounding it is the subject of §7.

(A note on the two parametrizations, so \texttt{λ} is unambiguous. The
action above uses the ADM/Hořava \texttt{λ} (the \texttt{K²}
coefficient); the mode-speed and PPN formulas use the Einstein-aether
couplings \texttt{c\_i}. On the relevant slice \texttt{c\_13\ =\ 0} they
are related by \texttt{λ\ =\ 1\ +\ c\_2}, so ED's
\texttt{λ\ ∈\ {[}0,\ 1/3)} (§7) is \texttt{c\_2\ ∈\ {[}−1,\ −2/3)}, and
the ADM healthy branch \texttt{λ\ \textless{}\ 1/3} coincides with the
aether no-ghost conditions on the low branch. Note this pins
\texttt{c\_2} to a narrow band, which constrains, and may be in tension
with, freely using the residual \texttt{c\_2} to set the MOND scale
\texttt{a\_0}; that compatibility is open.)

(In the equivalence-principle limit \texttt{η\ →\ 0}, i.e.~aether
coupling \texttt{c\_14\ →\ 0}, the standard scalar-mode speed formula
\texttt{c\_S²\ =\ c\_2(2\ −\ c\_14)/{[}c\_14(2\ +\ 3c\_2){]}} diverges,
so the scalar becomes effectively instantaneous, decoupling from
\emph{radiation} while surviving as a \emph{static} field. That is the
clean limit: near-pure-tensor, GR-like gravitational waves, plus a
static scalar in the sector where galactic MOND would live, with no
scalar-wave admixture to over-radiate. The \texttt{c\_14}-smallness is
itself open, and this limit is suggestive, not forced.)

\subsection{7. Computing λ: the Healthy Low Branch, λ ∈ {[}0, 1/3), Not
GR}\label{computing-ux3bb-the-healthy-low-branch-ux3bb-0-13-not-gr}

\texttt{λ} is the compression-vs-shear kinetic weighting, and it has an
exact handle. Decompose the extrinsic curvature into its trace and
traceless parts, \texttt{K\_ij\ =\ A\_ij\ +\ (1/3)\ h\_ij\ K}, where
\texttt{A\_ij} is the traceless \textbf{shear} (spin-2) and \texttt{K}
the trace \textbf{compression} (spin-0). Then, by exact algebra (this is
the orthogonality of the trace/traceless split, true for any symmetric
tensor; a probe confirms the identity numerically to machine precision,
\texttt{6×10⁻¹⁶}, which checks only the algebra, not anything
ED-specific),

\[ K_{ij}K^{ij} \;-\; \lambda K^2 \;=\; A_{ij}A^{ij} \;+\; \Big(\tfrac13 - \lambda\Big) K^2. \]

So the shear kinetic weight is \texttt{1} and the compression kinetic
weight is \texttt{w\_K\ =\ 1/3\ −\ λ}. \textbf{\texttt{λ} is nothing but
the compression-vs-shear weighting}, \texttt{λ\ =\ 1/3\ −\ w\_K}. This
gives three derived landmarks, each with a clear physical meaning:

{\def\LTcaptype{none} % do not increment counter
\begin{longtable}[]{@{}
  >{\raggedright\arraybackslash}p{(\linewidth - 4\tabcolsep) * \real{0.3333}}
  >{\raggedright\arraybackslash}p{(\linewidth - 4\tabcolsep) * \real{0.3333}}
  >{\raggedright\arraybackslash}p{(\linewidth - 4\tabcolsep) * \real{0.3333}}@{}}
\toprule\noalign{}
\begin{minipage}[b]{\linewidth}\raggedright
\texttt{λ}
\end{minipage} & \begin{minipage}[b]{\linewidth}\raggedright
compression weight \texttt{w\_K}
\end{minipage} & \begin{minipage}[b]{\linewidth}\raggedright
meaning
\end{minipage} \\
\midrule\noalign{}
\endhead
\bottomrule\noalign{}
\endlastfoot
\textbf{0} & \texttt{+1/3} & compression is a \textbf{normal
positive-energy mode}, weighted like the shear (the isotropic / identity
substrate kinetic term \texttt{K\_ij\ K\^{}ij}). Minimal. \\
\textbf{1/3} & \texttt{0} & compression carries \textbf{no} kinetic
energy (non-dynamical / fully constrained). \\
\textbf{1 (GR)} & \texttt{−2/3} & compression is a
\textbf{negative-energy conformal ghost} (the DeWitt supermetric). \\
\end{longtable}
}

The load-bearing fact about GR's value: \textbf{\texttt{λ\ =\ 1} is
\emph{forced} by full four-dimensional diffeomorphism invariance.} GR's
kinetic term is fixed by requiring the action be invariant under all
spacetime diffeomorphisms (the Einstein-Hilbert structure), and that
requirement makes the trace/conformal mode a negative-norm ghost,
rendered harmless by the Hamiltonian constraint. A \textbf{khronometric
substrate has only spatial diffeomorphisms plus the arrow's preferred
time}, not full 4D diffeomorphism invariance, so it is \emph{not} forced
to \texttt{λ\ =\ 1}. \texttt{λ} becomes a genuine free parameter set by
the substrate's inertia. (This is exactly why Hořava / khronometric
gravity carries \texttt{λ} as a coupling rather than fixing it to 1.)

ED's placement rests on the substrate inertia, and has three pieces at
three different tiers:

\begin{itemize}
\item
  \textbf{Derived (non-circular): ED is not \emph{forced} to
  \texttt{λ\ =\ 1}.} GR's value is forced by 4D diffeomorphism
  invariance (above); a khronometric substrate lacks that invariance, so
  \texttt{λ\ =\ 1} is not imposed and \texttt{λ} is a free parameter.
  This is clean and does not assume ED differs from GR; it only removes
  the constraint that would fix \texttt{λ\ =\ 1}.
\item
  \textbf{Natural value \texttt{λ\ =\ 0} (association-tier).} The
  kinetic term is supplied by ED's coherence / participation-velocity
  channel (the reversible participation ``velocity'' of the substrate;
  §8, an association). That inertia treats the compression and shear
  modes alike, giving each a positive kinetic weight with no built-in
  preference splitting trace from traceless. Equal positive weight is
  the identity supermetric, \texttt{K\_ij\ K\^{}ij}, which is
  \texttt{λ\ =\ 0}. So ED's \emph{natural / minimal} value is 0. This is
  a structural argument resting on the flagged Coh→kinetic association,
  not a computation.
\item
  \textbf{Upper limit \texttt{1/3} = the fully-constrained-compression
  endpoint; the interior is open.} Bandwidth conservation (P04) ties the
  compression mode \texttt{∂\_t\ b} to the current via continuity,
  \texttt{∂\_t\ b\ +\ ∇·J\ =\ 0}. In the limit where that continuity
  constraint \emph{fully} removes the compression mode's independent
  kinetic content, \texttt{w\_K\ →\ 0}, i.e.~\texttt{λ\ →\ 1/3} (the
  non-dynamical-compression endpoint). So conservation acts in the
  direction of suppressing the compression weight (0 → 1/3), and the
  fully-constrained endpoint is \texttt{1/3}. The claim that it can only
  ever suppress (never drive \texttt{w\_K} negative into GR's conformal
  ghost) is the \emph{positive-energy-substrate premise}, an assertion
  (ED's modes are positive-energy), not a theorem; it is what keeps ED
  off the ghost side. Where in the interior \texttt{{[}0,\ 1/3)} ED
  actually sits (how strongly conservation partially constrains the
  mode) is open.
\end{itemize}

Combining these:

\[ \boxed{\;\lambda \in [0,\,1/3),\quad \text{natural value } 0,\quad \lambda \neq 1\ \text{given ED's positive-energy substrate inertia}.\;} \]

Three consequences make this a physical result, not a formal one:

\begin{enumerate}
\def\labelenumi{\arabic{enumi}.}
\item
  \textbf{The tensor sector is GR-like for any \texttt{λ}.} The shear
  (spin-2) weight is \texttt{1} in ED and GR alike; \texttt{λ} never
  touches it. So ED's tensor graviton and its tensor radiation are
  GR-like \emph{regardless of where \texttt{λ} sits in
  \texttt{{[}0,\ 1/3)}}. This is why the pulsar and LIGO tensor matches
  survive the deformation.
\item
  \textbf{\texttt{{[}0,\ 1/3)} is the \emph{healthy} branch.} The
  standard khronometric stability result is that the scalar graviton is
  ghost-free and gradient-stable exactly on the low branch
  \texttt{λ\ \textless{}\ 1/3} (and on \texttt{λ\ \textgreater{}\ 1});
  the window \texttt{1/3\ \textless{}\ λ\ \textless{}\ 1} is unhealthy,
  and \texttt{λ\ =\ 1} is GR (scalar non-dynamical). ED's
  \texttt{{[}0,\ 1/3)} sits squarely on the healthy low branch, and
  \emph{not by tuning}: the positive-energy premise plus the
  conservation-suppression direction are what put it there. So ED's
  extra scalar (the khronon) is a well-behaved positive-energy field,
  which is the \emph{candidate} home for the MOND sector; that this
  scalar actually \emph{reproduces} MOND (rotation curves, baryonic
  Tully-Fisher) is not shown here and is inherited/open (§10.4).
\item
  \textbf{Not GR (given the premise).} \texttt{λ\ ≠\ 1} (on the
  positive-energy premise) means ED's metric kinetic structure differs
  from GR's in the trace sector, on the healthy side of it. ED is a
  genuine deformation of GR, not GR. The clean, unconditional part is
  weaker but solid: ED is \emph{not forced} to GR's \texttt{λ\ =\ 1}, so
  nothing in the khronometric structure requires it to be GR.
\end{enumerate}

\textbf{Honest tier (preamble 3).} The decomposition and the three
landmarks are derived. ``\texttt{λ\ =\ 0} natural, bounded to
\texttt{{[}0,\ 1/3)}, healthy branch'' is a \emph{structural argument}
resting on the coherence-to-kinetic association (§8), not a computed
decimal. \texttt{λ\ ≠\ 1} and the healthy-branch placement are robust
given positive-energy substrate modes. The residual, exactly
\emph{where} in \texttt{{[}0,\ 1/3)} (which fixes the precise
scalar-mode speed and the MOND interpolation constant), is how strongly
conservation suppresses the compression weight, and needs the
mode-resolved P04+P05 inertia. That is the one calculation left in ED
gravity.

\subsection{8. The Substrate Map Behind the Bound (Honest
Provenance)}\label{the-substrate-map-behind-the-bound-honest-provenance}

The \texttt{λ} bound rests on a map from the P12 stability landscape
\texttt{Σ\ =\ Coh\ −\ Str\ −\ Grad} to the effective action, and honesty
requires marking which parts of that map are verified and which are
association.

Canonical P12 has \texttt{a\ =\ −∇Σ} (Newton's second law), so
\textbf{\texttt{Σ} is a potential}, and the \emph{kinetic} term (the
inertia that makes the geometry radiative) is a separate structure.
Organized by the author's ``Relation-Boundary-Gradient'' reading of
\texttt{Coh\ −\ Str\ −\ Grad} and an older, non-canonical ED cosmology
manuscript (flagged as an association, not a canonical source, per
preamble 4), the map is three-way:

{\def\LTcaptype{none} % do not increment counter
\begin{longtable}[]{@{}
  >{\raggedright\arraybackslash}p{(\linewidth - 6\tabcolsep) * \real{0.2500}}
  >{\raggedright\arraybackslash}p{(\linewidth - 6\tabcolsep) * \real{0.2500}}
  >{\raggedright\arraybackslash}p{(\linewidth - 6\tabcolsep) * \real{0.2500}}
  >{\raggedright\arraybackslash}p{(\linewidth - 6\tabcolsep) * \real{0.2500}}@{}}
\toprule\noalign{}
\begin{minipage}[b]{\linewidth}\raggedright
reading
\end{minipage} & \begin{minipage}[b]{\linewidth}\raggedright
\texttt{Σ} term
\end{minipage} & \begin{minipage}[b]{\linewidth}\raggedright
effective-action role
\end{minipage} & \begin{minipage}[b]{\linewidth}\raggedright
status
\end{minipage} \\
\midrule\noalign{}
\endhead
\bottomrule\noalign{}
\endlastfoot
\textbf{Gradient} & Grad & spatial curvature \texttt{ξ\ ³R} &
\textbf{verified} \\
\textbf{Boundary} & Str & potential / mass term & association \\
\textbf{Relation} & Coh & the \textbf{kinetic} term
\texttt{K\_ij\ K\^{}ij\ −\ λK²} (hence \texttt{λ}) & association \\
\end{longtable}
}

The \textbf{verified} piece: the emergent spatial metric is conformal,
\texttt{h\_ij\ =\ (1/b)δ\_ij}, and its scalar curvature (probe-checked
to finite-difference accuracy) is

\[ {}^3R \;=\; 2\nabla^2 b \;-\; \tfrac52\,\frac{(\nabla b)^2}{b}, \]

whose \texttt{(\textbackslash{}nabla\ b)\^{}2/b} piece is exactly
\texttt{Σ}'s \textbf{Grad} term (the \texttt{∇²b} piece is a boundary
total-derivative). So the spatial-curvature term of the action \emph{is}
the substrate's gradient penalty. The gradient sector of the map is
established.

The \textbf{association}: identifying the \emph{kinetic} term (hence
\texttt{λ}) with \texttt{Σ}'s coherence sector is a pointer from the RBG
reading and the older manuscript, not a canonical derivation. Its value
is that it locates \emph{which} part of the substrate supplies the
inertia (the reversible participation channel), which is what sharpens
\texttt{λ} to the single compression-vs-shear weighting of §7. It also
matches that manuscript's finding that the inertial/wave sector
dominates near the density ceiling (high acceleration), consistent with
``high acceleration → GR/wave regime'' (pulsars). But it is an
association: the canonical grounding of the coherence kinetic term, and
with it the exact point of \texttt{λ} in \texttt{{[}0,\ 1/3)}, is the
open work.

\subsection{9. Limitations (honest)}\label{limitations-honest}

\begin{itemize}
\item
  \textbf{Not GR; the deformation parameter is bounded, not pinned.} ED
  is khronometric, not GR (preamble 1). \texttt{λ\ ∈\ {[}0,\ 1/3)} is a
  signed bound with a natural value 0, not a decimal (§7); the exact
  point needs the mode-resolved conservation inertia and is open.
\item
  \textbf{The khronometric identification is account-tier.} ``Arrow →
  preferred frame → khronometric class'' is a structural identification
  with a standard viable class (preamble 2, §3), not a bespoke
  derivation of every coupling. \texttt{ξ\ =\ 1} and \texttt{η\ →\ 0}
  are structural conditions (§5-6), not first-principles coupling
  computations; their exact values, and the full BBN /
  gravitational-Cherenkov fit, need the coarse-grained effective action,
  which is not built.
\item
  \textbf{The coherence-to-kinetic map is an association.} The kinetic
  sector's identification with \texttt{Σ}'s coherence term (§8) is
  organized by a non-canonical manuscript and the RBG reading. Only the
  gradient sector (\texttt{Grad\ →\ ³R}) is verified. Treat the older
  cosmology manuscript as an association, not a canonical source.
\item
  \textbf{\texttt{c\_T\ =\ c} carries a load-bearing assumption.} It is
  structural and confirmed by GW170817 (§5), but it assumes ED's
  universality holds to 10⁻¹⁵ through the radiative sector (preamble 5);
  a radiative detuning of the light cone is not excluded.
\item
  \textbf{Values inherited.} \texttt{G}, \texttt{ℓ\_P}, \texttt{a₀} are
  not derived here.
\item
  \textbf{Shares the companion paper's layer caveat.} The emergent
  metric and its spin-2 channel are established in the counting /
  coarse-grained (layer-2) description; that ED's raw ballistic
  (layer-1) dynamics coarse-grains to them is attributed, not shown (see
  the companion paper's §7).
\end{itemize}

\subsection{10. Falsification Criteria}\label{falsification-criteria}

\subsubsection{10.1 ED gravity is not khronometric (wrong field
content)}\label{ed-gravity-is-not-khronometric-wrong-field-content}

\textbf{Falsifier sentence:} \emph{A demonstration that ED's emergent
gravity does not carry a spin-2 tensor graviton (that the emergent
metric has no genuine traceless dynamical channel for a non-spherical
source, contrary to §4), or that ED has no preferred frame (contrary to
the arrow being a genuine physical foliation), or that ED's khronon is
purely rigid (the emergent, b-sourced foliation fluctuation is not in
fact dynamical, contrary to the §3 structural assumption, which would
remove the scalar mode and the MOND sector entirely), would falsify the
khronometric identification and return ED to either a falsified scalar
theory or a scalar-less rigid-background theory.}

\subsubsection{10.2 The light cone is not
shared}\label{the-light-cone-is-not-shared}

\textbf{Falsifier sentence:} \emph{A gravitational-wave speed differing
from \texttt{c} beyond the GW170817 bound
(\texttt{\textbar{}c\_T\ −\ c\textbar{}/c\ \textless{}\ 10⁻¹⁵}), or a
detected scalar-mode (breathing) polarization admixture above LIGO's
bounds, would falsify ED's structural \texttt{c\_T\ =\ c} and the
equivalence-principle suppression of the scalar radiative sector (§5).}

\subsubsection{10.3 λ is outside the healthy branch, or is
1}\label{ux3bb-is-outside-the-healthy-branch-or-is-1}

\textbf{Falsifier sentence:} \emph{A computation of ED's
compression-vs-shear kinetic weighting yielding \texttt{λ\ ≥\ 1/3} (the
unhealthy window or beyond) would falsify the healthy-low-branch result
(§7); yielding \texttt{λ\ =\ 1} would collapse ED to GR and contradict
the khronometric structure (and the observed MOND phenomenology); a
negative compression weight (\texttt{λ\ \textgreater{}\ 1/3} toward a
ghost) would falsify the positive-energy-substrate premise.}

\subsubsection{10.4 The scalar sector does not give
MOND}\label{the-scalar-sector-does-not-give-mond}

\textbf{Falsifier sentence:} \emph{If the healthy scalar (khronon)
sector, in its low-acceleration limit, does not reproduce the MOND
phenomenology (galactic rotation curves and the baryonic Tully-Fisher
relation) but instead requires particle dark matter, the identification
of ED's \texttt{λ}-sector scalar with the MOND sector (§6-7) is
falsified.}

\subsubsection{10.5 Pulsar decay inconsistent with GR-tensor-dominant
radiation}\label{pulsar-decay-inconsistent-with-gr-tensor-dominant-radiation}

\textbf{Falsifier sentence:} \emph{A binary-pulsar orbital-decay rate
inconsistent with a GR-tensor-dominant luminosity plus the
equivalence-principle-suppressed small scalar-quadrupole correction (an
over- or under-radiation outside the viable khronometric window) would
falsify ED's placement in that window (§5).}

\subsection{11. Position Statement}\label{position-statement}

\textbf{ED gravity is a healthy one-parameter khronometric deformation
of General Relativity.} ED has a preferred frame, the arrow (P11), so
its emergent gravity is Lorentz-violating and falls in the
Einstein-aether / Hořava / khronometric class with the aether equal to
the arrow. That class carries GR's spin-2 tensor graviton (the emergent
metric develops a genuine traceless part for a non-spherical source,
kinematically present in a static probe; its propagation as a graviton
follows from the class), so its tensor radiation is GR-like and matches
binary pulsars and LIGO. ED's minimal, universal setup selects the
viable corner via three \emph{distinct} conditions, not one mechanism:
its universal coupling (equivalence principle) structurally evades the
scalar dipole catastrophe (the one genuinely ED-structural leg), while a
shared graviton/photon light cone (\texttt{c\_13\ =\ 0},
\texttt{c\_T\ =\ c}, a confirmed consistency against GW170817's 10⁻¹⁵
bound) and small preferred-frame PPN couplings (\texttt{c\_14\ →\ 0},
open) are separate structural conditions its minimality makes natural
but does not derive. The effective action is the khronometric form with
\texttt{ξ\ =\ 1} (light cone) and \texttt{η\ →\ 0} (equivalence
principle), leaving \texttt{λ} (the \texttt{K²} coefficient) as the
single distinctive parameter. Bounding \texttt{λ} as the
compression-vs-shear kinetic weighting places it on the healthy low
branch, \texttt{λ\ ∈\ {[}0,\ 1/3)}, natural value 0, \texttt{λ\ ≠\ 1}
given ED's positive-energy substrate inertia (the clean unconditional
part being that ED is \emph{not forced} to GR's \texttt{λ\ =\ 1}, which
a 4D diffeomorphism invariance ED's substrate lacks would impose), on
which the extra scalar is a well-behaved positive-energy field, the
candidate home for galactic MOND (that it reproduces MOND is
inherited/open).

\textbf{What this paper claims.} The khronometric identification
(grounded in the arrow), the spin-2 channel kinematically present
(static probe; propagation from the class), the radiative viability with
the dipole evaded (ED-structural) and \texttt{c\_T\ =\ c} a confirmed
consistency (a structural condition its minimality makes natural), the
effective-action form and its reduction to one parameter (grounded), and
the bound \texttt{λ\ ∈\ {[}0,\ 1/3)} on the healthy low branch (derived
landmarks + structural placement on the positive-energy premise).
Together: \textbf{ED gravity is GR at high acceleration and MOND at low
acceleration, a healthy Lorentz-violating deformation of GR with a
single bounded distinctive parameter.}

\textbf{What this paper does not claim.} Exact GR (the Einstein equation
is not derived); a computed decimal for \texttt{λ} (bounded, not pinned;
the exact point is open); first-principles khronometric couplings
(\texttt{ξ\ =\ 1}, \texttt{η\ →\ 0} are structural conditions, not
derived numbers; the coarse-grained effective action is unbuilt); a
canonical coherence-to-kinetic map (an association; only
\texttt{Grad\ →\ ³R} is verified); derived values (\texttt{G},
\texttt{ℓ\_P}, \texttt{a₀} inherited); or that ED's raw ballistic
dynamics reaches these (a layer-2 attribution).

\textbf{Series context.} The dynamical companion to the kinematics paper
(Relational Gravity: the emergent metric
\texttt{g\ \textasciitilde{}\ 1/b}, Gauss's law, relational character,
MOND nonlinearity) and to Paper\_027 (Newton's \texttt{G}). Where the
kinematics paper fixes the geometry and Paper\_027 fixes the
coefficient, this paper fixes the \emph{dynamical class} (khronometric),
its radiative viability (\texttt{c\_T\ =\ c}), and its one distinctive
parameter (\texttt{λ\ ∈\ {[}0,\ 1/3)}). The preferred frame is the arrow
the reduction program isolates as the keystone primitive; ED gravity is
the shadow that frame casts on the emergent metric.

\begin{center}\rule{0.5\linewidth}{0.5pt}\end{center}

\subsection{Appendix: Glossary and Reader
Map}\label{appendix-glossary-and-reader-map}

\subsubsection{Glossary}\label{glossary}

\begin{itemize}
\tightlist
\item
  \textbf{Khronometric gravity.} The low-energy effective theory of a
  preferred-frame (Lorentz-violating) gravity, with a global time
  function (the ``khronon'') defining a preferred foliation; the
  hypersurface-orthogonal limit of Einstein-aether theory, and the
  low-energy limit of Hořava gravity. ED's preferred frame is the arrow
  (P11).
\item
  \textbf{Aether / preferred frame.} A dynamical unit timelike vector
  \texttt{u\^{}μ} picking out ``space at a time.'' In ED,
  \texttt{u\^{}μ} = the arrow, with the §3 caveat that the
  \emph{primitive} arrow fixes the foliation's existence/orientation
  (rigid) while the \emph{dynamical} khronon that propagates the scalar
  mode is the emergent, \texttt{b}-sourced fluctuation.
\item
  \textbf{Spin-2 tensor graviton.} The two transverse-traceless
  polarizations of GR's gravitational wave. The khronometric class
  carries them, so its tensor radiation is GR-like. ED's emergent metric
  develops the required traceless channel for a non-spherical source
  (§4).
\item
  \textbf{Scalar (khronon) mode.} The one extra propagating mode of
  khronometric gravity beyond GR's tensor graviton. In ED it is the
  candidate MOND sector (that it reproduces MOND is inherited/open), and
  \texttt{λ\ ∈\ {[}0,\ 1/3)} makes it healthy (positive-energy).
\item
  \textbf{\texttt{λ} (the compression-vs-shear weighting).} The
  coefficient of \texttt{K²} in the kinetic term; equals
  \texttt{1/3\ −\ w\_K} where \texttt{w\_K} is the compression (trace)
  mode's kinetic weight relative to the shear. GR is \texttt{λ\ =\ 1} (a
  conformal ghost, forced by 4D diffeomorphism invariance); ED is
  \texttt{λ\ ∈\ {[}0,\ 1/3)} (healthy, natural value 0).
\item
  \textbf{\texttt{c\_T\ =\ c}.} Gravitational-wave speed equal to light
  speed; a gravity-sector condition (\texttt{c\_13\ =\ 0}, a shared
  graviton/photon cone) that ED's minimal setup makes natural, confirmed
  by GW170817 to 10⁻¹⁵. A consistency, not a coupling ED derives.
\item
  \textbf{Universal coupling / equivalence principle.} Mass = universal
  bandwidth-influence \texttt{Q\ ∝\ M}; the emergent metric is one
  universal read-out of \texttt{b}, so all bodies fall alike. This
  \emph{matter-sector} feature is what structurally evades the scalar
  dipole. It does \emph{not} by itself give \texttt{c\_T\ =\ c} or small
  PPN (those are separate gravity-sector conditions,
  \texttt{c\_13\ =\ 0} and \texttt{c\_14\ →\ 0}); ED's minimality makes
  all three natural, but only the dipole evasion is a direct consequence
  of the universal coupling.
\item
  \textbf{Healthy low branch.} The khronometric stability window
  \texttt{λ\ \textless{}\ 1/3}, on which the scalar mode is ghost-free
  and gradient-stable. ED lands here structurally.
\end{itemize}

\subsubsection{Reader map}\label{reader-map}

\emph{Intuition.} ED has an arrow of time that is real, not emergent, so
ED's spacetime has a preferred ``now.'' A gravity built on such a
spacetime cannot be exactly Einstein's (which has no preferred now); it
is the next thing over, a preferred-frame gravity, and physicists have a
well-studied name and toolkit for that: khronometric / Einstein-aether
gravity. The good news is that this kind of gravity still has ordinary
gravitational waves (the spin-2 graviton), so it passes the pulsar and
LIGO tests, and it comes with one extra ingredient, a scalar mode, which
is exactly where galactic MOND can live. ED is a \emph{minimal,
universal} version of such a gravity, and that minimality makes it land
in the small ``viable corner'' that observations allow. Three separate
things have to hold in that corner, and it is worth being precise that
they are separate: everything falling the same way (ED's equivalence
principle) directly removes the dangerous extra radiation; gravitational
waves travelling at the speed of light (which the 2017 neutron-star
merger confirmed) is a distinct condition on the gravity sector that
ED's minimality makes natural; and keeping the preferred-frame effects
tiny is a third such condition. ED does not derive all three from one
knob; its minimality makes the corner where all three hold the natural
one. The whole theory then comes down to a single number, \texttt{λ},
measuring how much the ``squeezing'' motion of space weighs against the
``shearing'' motion. GR is \emph{forced} to one value of that number by
a symmetry ED does not have; ED is therefore not forced to Einstein's
value, and its own inertia and conservation law point to a small healthy
range instead, \texttt{0\ ≤\ λ\ \textless{}\ 1/3}, on the well-behaved
side (this last placement resting on how we read ED's inertia, an
association, not a bare calculation).

\emph{The distinctive claim.} ED gravity is a healthy one-parameter
khronometric deformation of GR: GR-like tensor waves, a shared light
cone, and a healthy scalar in the sector where MOND would live
(reproduction inherited/open), with its one free parameter bounded to
\texttt{{[}0,\ 1/3)}. It is GR at high acceleration and MOND at low,
from one preferred-frame structure.

\textbf{Where to look next.} - \emph{The emergent metric this puts
dynamics on:} the kinematics paper (Relational Gravity,
\texttt{g\ \textasciitilde{}\ 1/b}). - \emph{The coefficient
\texttt{G}:} Paper\_027 (Newton's G). - \emph{The interference / MOND
mechanism in the nonlinear sector:} the P14 / interference-gravity
result. - \emph{The preferred frame as the keystone primitive:} the
arrow / minimal-ontology work.

\textbf{Open work} (declared): the exact point of \texttt{λ} in
\texttt{{[}0,\ 1/3)} (needs the mode-resolved conservation inertia;
fixes the precise scalar-mode speed and the MOND interpolation
constant); the coarse-grained effective action that would compute ED's
khronometric couplings from first principles (rather than fixing them by
the structural conditions \texttt{ξ\ =\ 1}, \texttt{η\ →\ 0}); and the
canonical grounding of the coherence kinetic term (currently an
association). The class, the radiative viability, and the one-parameter
bound are established; the decimal is the frontier.

\begin{center}\rule{0.5\linewidth}{0.5pt}\end{center}

\textbf{End of Paper (Khronometric Gravity).}

\emph{Gravity arc, dynamical companion. ED has a preferred frame (the
arrow, P11), so its emergent gravity is Lorentz-violating: the
khronometric / Einstein-aether / Hořava class, aether = arrow. This
carries GR's spin-2 tensor graviton (the emergent metric develops a
genuine traceless part for a non-spherical source, kinematically present
in a static probe; propagation from the class), so tensor radiation is
GR-like (pulsars, LIGO). ED's minimal, universal setup selects the
viable corner via three distinct conditions, not one mechanism: its
universal coupling (equivalence principle, \texttt{Q\ ∝\ M})
structurally evades the scalar dipole (the one ED-structural leg), while
a shared graviton/photon light cone (\texttt{c\_13\ =\ 0},
\texttt{c\_T\ =\ c}, a confirmed consistency with GW170817) and small
preferred-frame PPN (\texttt{c\_14\ →\ 0}, open) are separate structural
conditions its minimality makes natural but does not derive. The
effective action is khronometric with \texttt{ξ\ =\ 1} and
\texttt{η\ →\ 0}, leaving \texttt{λ} as the one distinctive parameter.
Bounding \texttt{λ} (the compression-vs-shear kinetic weighting) places
it on the healthy low branch \texttt{{[}0,\ 1/3)}, natural value 0,
\texttt{λ\ ≠\ 1} given ED's positive-energy substrate inertia (the clean
part being that ED is not }forced* to GR's \texttt{λ\ =\ 1}, which a 4D
diffeomorphism invariance ED lacks would impose). So ED gravity is a
healthy one-parameter khronometric deformation of GR: GR at high
acceleration, and at low acceleration a healthy scalar in the sector
where MOND would live (reproduction inherited/open). Open: the exact
point in \texttt{{[}0,\ 1/3)}. Companion to the kinematics paper and
Paper\_027.*

\end{document}
